% Document class and font size
\documentclass[a4paper,9pt]{extarticle}

% Packages
\usepackage[utf8]{inputenc}
\usepackage{geometry}
\geometry{letterpaper, margin=0.75in}
\usepackage{titlesec}
\usepackage{enumitem}
\usepackage{hyperref}
\usepackage{tabularx}
\usepackage{fancyhdr}
\usepackage[backend=bibtex,style=ieee]{biblatex}
\bibliography{ref}

% Formatting
\setlist{noitemsep}
\titleformat{\section}{\large\bfseries}{\thesection}{1em}{}[\titlerule]
\titlespacing*{\section}{0pt}{\baselineskip}{\baselineskip}

\newcommand\mycite[1]{\AtNextCite{\defcounter{maxnames}{1}}\citeauthor{#1}, \citeyear{#1}}

% Begin document
\begin{document}

% Disable page numbers
\pagestyle{fancy}
\renewcommand{\headrulewidth}{0pt}
\fancyhead{}
\fancyhead[L]{\textit{Aditya Dendukuri}}
\fancyhead[R]{\textit{PhD Candidate at UCSB}}
\thispagestyle{empty}

% Header
\begin{center}
    \textbf{\LARGE Aditya Dendukuri}\\[2pt]
    PhD Candidate \\ University of California, Santa Barbara\\[6pt]
    \href{mailto:dendukuri.aditya123@outlook.com}{dendukuri.aditya123@outlook.com} \\ 
    \href{https://adityadendukuri.github.io/}{adityadendukuri.github.io} 
\end{center}

% Research Interests Section
\section*{Research Interests}
\noindent
Numerical analysis, mathematical modeling and scientific machine learning. 

% Education Section
\section*{Education}
\noindent
\textbf{University of California, Santa Barbara CA}\hfill Sept 2021 | Present\\
PhD in Computer Science \hfill \\
Advisor: Linda R. Petzold 
%\textit{Expected Graduation: Winter 2027}

\noindent\\
\textbf{University of Arkansas, Fayetteville AR} \hfill Aug 2018 | Jul 2020\\
MS in Computer Science \\
Advisors: David M. Ford and Matthew J. Patitz \\
Thesis: Nonlinear Dimensionality Reduction for the Thermodynamics of Small Clusters of Particles

\noindent\\
\textbf{University of Central Arkansas, Conway AR} \hfill Aug 2014 | May 2018\\
BS in Computer Science (Minor in Physics)\\
Advisor: Tansel Halic \\



\section*{Experience}

\noindent
\textbf{Lawrence Berkeley National Laboratory} \hfill Berkeley, CA\\
\textit{UC Graduate Student Researcher (GSR) Intern} \hfill July 2026 -- present
\begin{itemize}
    \item Develop efficient approximation methods for stochastic processes, particularly continuous-time Markov chains arising in chemical kinetics and reaction networks.
    \item Devise mathematical methods and implement computational programs for accurately simulating stochastic processes arising in physical science and artificial intelligence applications.
    \item Collaborate with scientists in the Applied Mathematics and Computational Research Division on computational methods in advanced mathematics and machine learning.
\end{itemize}

\noindent
\textbf{NSF National Center for Atmospheric Research (NCAR)} \hfill Boulder, CO\\
\textit{SIParCS Graduate Student Intern} \hfill May 2024 -- July 2024
\begin{itemize}
    \item Contributed to translating a radiative transfer codebase from Fortran to C++ and prepared documentation to support maintainability and future development.
    \item Implemented a custom solver for tridiagonal linear systems and benchmarked it against LAPACKE to evaluate numerical accuracy and performance.
    \item Extended testing coverage for critical linear algebra components using Google Test and added performance benchmarks with Google Benchmark.
\end{itemize}

\noindent
\textbf{University of Arkansas} \hfill Fayetteville, AR\\
\textit{Research Associate} \hfill Sept 2020 -- June 2021
\begin{itemize}
    \item Modeled nucleation and phase-transition processes using Monte Carlo simulations and implemented custom C code to analyze critical-cluster formation dynamics.
\end{itemize}

\section*{Research Assistantships}

\noindent
\textbf{University of California, Santa Barbara} \hfill Santa Barbara, CA\\
\textit{Graduate Research Assistant} \hfill Sept 2021 -- Present
\begin{itemize}
    \item Develop structure-exploiting numerical algorithms for analyzing high-dimensional discrete stochastic systems, enabling efficient simulation and inference for complex reaction networks.
    \item Investigate surrogate modeling approaches for understanding activity in lab-grown neural tissues, using time-series analysis and dimensionality reduction to extract interpretable spatiotemporal dynamics.
\end{itemize}

\noindent
\textbf{University of Arkansas} \hfill Fayetteville, AR\\
\textit{Graduate Research Assistant} \hfill Aug 2018 -- July 2020
\begin{itemize}
    \item Explored the thermodynamic properties of colloidal clusters using data-driven methods, reconstructing energy landscapes and transition pathways from stochastic simulations.
    \item Simulated quantum circuits to investigate hybrid quantum-classical learning models and developed workflows for evaluating variational quantum algorithms.
    \item Published results on data-driven thermodynamics in \textit{Computers \& Chemical Engineering}.
\end{itemize}

\noindent
\textbf{University of Central Arkansas} \hfill Conway, AR\\
\textit{Undergraduate Research Assistant} \hfill Aug 2016 -- May 2018
\begin{itemize}
    \item Developed real-time surgical training simulations by coupling fluid-flow and heat-transfer models in a numerical solver.
    \item Conducted hierarchical task analysis using observational data from arthroscopic surgeries to inform system design and evaluation.
    \item Presented and published this work at the 23rd IEEE International Conference on Bioinformatics and Bioengineering.
\end{itemize}

\section*{Teaching Assistantships}

\noindent
\textbf{University of California, Santa Barbara} \hfill Santa Barbara, CA\\
\textit{Graduate Teaching Assistant} \hfill Sept 2021 -- Present
\begin{itemize}
    \item Supported undergraduate instruction in scientific computing, data structures, and numerical methods, including programming in Python, C++, and MATLAB.
\end{itemize}
\section*{Preprints (Under Review/ In Press)}
\noindent
\begin{itemize}
    \item \fullcite{dendukuri2026adjointkrylov}
    \item \fullcite{dendukuri2026koopmancme}
    \item \fullcite{dendukuri2025fluxpreservingadaptivefinitestate} 
    \begin{itemize}
        \item Accepted for publication (in press) at \textit{SIAM Multiscale Modeling and Simulation}.
    \end{itemize}
\end{itemize}

% Publications Section
\section*{Publications}
\noindent
\begin{itemize}
    \item \fullcite{van_der_molen_preconfigured_2026}    
    \item \fullcite{dendukuri:bibe}
    \item \fullcite{ford2020machine}
    \item \fullcite{demirel2017hierarchical}
\end{itemize}

% Conferences Section
\section*{Conference Presentations}
\begin{itemize}
    \item \href{https://meetings.siam.org/sess/dsp_programsess.cfm?SESSIONCODE=89397}{\textbf{Data-Driven Rate Recovery and Sensitivity Analysis for Discrete Stochastic Chemical Reaction Systems.}} SIAM Annual Meeting, Cleveland, OH, 2026.
    \item \fullcite{rdme_boulder}
    \item \fullcite{dendukurip:bibe}
    \item \fullcite{siam}
    \item \fullcite{fopam}
    \item \fullcite{qtml}
    \item \fullcite{inbre}
    \item \fullcite{aas}
\end{itemize}

\section*{Professional Service}
\begin{itemize}
    \item Session chair, \href{https://meetings.siam.org/sess/dsp_programsess.cfm?SESSIONCODE=89397}{CP17: Stochastic Systems}, SIAM Annual Meeting, Cleveland, OH, 2026.
\end{itemize}

% Honors Section
\section*{Honors}
\begin{itemize}[leftmargin=*]
    \item Best paper award. IEEE Bioinformatics and Bioengineering (BIBE) Conference 2023. 
    \item Masters thesis nominated for the \textit{Conference of Southern Graduate Schools (CSGS)} outstanding master's thesis award.
    \item Graduate Student Council Travel Grant (\$500) for the Foundations of Process Analytics and Machine Learning Conference in Raleigh, North Carolina (Summer 2019).
    \item Advancement of Undergraduate Research in the Sciences (AURS) summer research fund (2016) (\url{http://www.aursfund.org}) Award Amount: \$5000
\end{itemize}

% Selected Courses Section
\section*{Selected Courses}
\begin{tabularx}{1\textwidth}{
>{\raggedright\arraybackslash}X
>{\raggedright\arraybackslash}X}
\textbf{PhD Courses (UCSB)}& \textbf{MS Courses (UofA)}\\
\begin{itemize}
    \item Matrix Analysis and Computation
    \item Level Set Methods
    \item Finite Element Methods
    \item Numerical Simulation
    \item Spectral Graph Theory
\end{itemize}
& \begin{itemize}
    \item Deep Learning
    \item Natural Language Processing
    \item Computer Vision
\end{itemize}\\
\end{tabularx}

% Skills Section
\section*{Skills (2-3 Years of Full-Time Experience)}
\begin{itemize}
    \item \textbf{Programming:} Julia, Python, C/C++
    \item \textbf{Software:} Tensorflow and PyTorch (Python), Modelling Toolkit (Julia), SciML (Julia), Linux, CMAKE.
\end{itemize}

% References Section
\section*{References} 
\textbf{}
\textbf{Dr. Linda R. Petzold} (petzold@engineering.ucsb.edu)\\
Mehrabian Distinguished Professor, Mechanical Engineering and Computer Science, University of California, Santa Barbara.  
\vspace{4pt}\\
\textbf{Jian Sun} (sunjian@ucar.edu)\\
Soft Eng/Prog II, NSF National Center for Atmospheric Research
\vspace{4pt}\\
\textbf{Dr. David Ford} (ford3dm@cmich.edu) \\
Dean, Bagley College of Engineering, Mississippi State University.\\
Past: Professor and Department Head. Ralph E. Martin Department of Chemical Engineering, University of Arkansas. 
\vspace{4pt}\\
\textbf{Dr. Tansel Halic} (tanselhalic@gmail.com) \\
Senior Software Engineer, Intuitive Surgical.

\end{document}
